\section{Разработка программного обеспечения интеллектуальной системы}
\subsection{Техническое задание}                   %           2. Тех задание
\subsubsection{Основание для разработки}  

Основанием для разработки является задание на выпускную квалификационную работу магистра «Интеллектуальная система оптимизации производственной деятельности на основе АИС» согласно приказу №1586-с от 03.04.2026 «Об утверждении тем выпускных квалификационных работ и руководителей выпускных квалификационных работ».

\subsubsection{Назначение разработки}

Программный продукт предназначен для автоматизации процессов производственного планирования и оптимизации диспетчеризации на предприятии. Функциональное назначение системы заключается в формировании выполнимых производственных расписаний, сокращении простоев оборудования и снижении времени опоздания заказов за счёт интеллектуального распределения операций по доступным ресурсам с учётом их текущего состояния, специализации и загрузки.

Система ориентирована на применение в цехах с гибкой (Flexible Job Shop) структурой, где одна операция может быть выполнена на нескольких альтернативных станках. Программный продукт может использоваться как самостоятельное решение для оперативного управления производством, так и в составе корпоративной информационной инфраструктуры, принимая заказы из ERP-систем и передавая события исполнения в MES-среду.

Дополнительной особенностью программного продукта является наличие встроенного модуля обучения и применения нейросетевых моделей диспетчеризации. Поддерживаются три парадигмы обучения \cite{prinreshen}: поведенческое клонирование (Behavioral Cloning, PointerNet), самоконтролируемое обучение с саморазметкой (Self‑Labeling, SLIM) и обучение с подкреплением (PPO). Модели способны прогнозировать оптимальный ресурс для каждой операции на основе накопленных производственных данных, учитывая такие факторы, как тип операции, нормативную длительность, приоритет заказа, оставшееся время до срока сдачи, текущую загрузку и коэффициент готовности ресурса. Это позволяет не только строить расписания в реальном времени, но и адаптироваться к изменениям производственной обстановки (выход станка в ремонт, изменение надёжности).

\subsubsection{Требования к программе}

\paragraph{Требования к данным программы}

Входными данными для интеллектуальной системы производственного планирования являются \cite{setiavtomat}:
\begin{itemize}
\item заказы на производство — перечень работ с указанием номенклатуры, приоритета, срока сдачи и состава технологических операций;
\item технологические операции — описание единичной работы: тип (фрезеровка, сборка, сварка), нормативная длительность, количество изготавливаемых деталей, предшественники (при необходимости);
\item производственные ресурсы — информация о станках и рабочих центрах: идентификатор, название, текущий статус (работает, авария, в ремонте), коэффициент готовности (надёжности), флаг нахождения в ремонте, допустимые типы операций (связи ресурс–тип);
\item данные о состоянии ресурсов — актуальные показатели загрузки (число ожидающих операций, наличие активной задачи), коэффициент готовности, признак ремонта, средняя загрузка за период;
\item обучающие выборки — синтетические или исторические данные о назначениях операций на ресурсы, сформированные по эвристическим правилам или MILP-решениям, включающие векторы признаков операций и ресурсов и целевые метки назначенных ресурсов.
Выходными данными системы являются:
\item производственное расписание — распределение операций по ресурсам с указанием времени старта и окончания каждой операции, сформированное на основе выбранного метода планирования (эвристика, нейросеть BC, SLIM или PPO);
\item рекомендации по назначению операций — результат работы нейросетевого агента, выдающего для каждой неразмещённой операции предпочтительный ресурс с учётом текущей производственной обстановки;
\item обученные нейросетевые модели — файлы весов для каждого из трёх методов обучения (PointerNet, SLIM, PPO), пригодные для последующей загрузки и использования в планировании без повторного обучения;
\item метрики обучения и сравнения — значения функции потерь, точности классификации на валидационной выборке, матрицы ошибок, отчёты о классификации, а также сводные KPI расписаний (среднее и максимальное опоздание, процент опаздывающих заказов, makespan, коэффициент вариации загрузки ресурсов) для каждого метода планирования; 
\item визуализация — диаграмма Ганта производственного плана, гистограммы опозданий, графики загрузки ресурсов, кумулятивные кривые выполнения, box‑plot и другие графические материалы.
\end{itemize}

\paragraph{Концептуальная модель предметной области}

Концептуальная модель системы представляет собой взаимосвязь основных сущностей и потоков данных.

Основные компоненты модели:
\begin{enumerate}
\item «Ядро» — основной модуль, где происходят вычисления, сортировка и подключение к остальным модулям.
\item «Файловое хранилище» — локальное хранилище пользователя, где установлена программа, все модели, метрические данные и дополнительные компоненты хранятся.
\item «База данных»  — централизованное хранилище заказов, операций, ресурсов, связей типов, очередей задач и результатов планирования.
\item «Интерфейс» — графическое приложение на PySide6, обеспечивающее визуализацию дашборда, диаграммы Ганта, редактирование заказов и ресурсов, запуск планирования и обучения.
\end{enumerate}

На рисунке ~\ref{koncept_model:image} представлена структурная схема концептуальной модели предметной области.

\begin{figure}[H]
\includegraphics[width=1\linewidth]{koncept_model}
\caption{Структурная схема концептуальной модели предметной области}
\label{koncept_model:image}
\end{figure}

Взаимодействие компонентов:
\begin{enumerate}
\item Пользователь через интерфейс создаёт заказы, настраивает ресурсы и их готовность.
\item Планировщик получает актуальные данные из базы, выбирает метод (эвристика или загруженная модель) и строит расписание.
\item При использовании нейросети обученные веса подгружаются из хранилища моделей, а признаки операции и ресурсов подаются на вход сети для получения рекомендованного ресурса.
\item Результаты планирования отображаются на дашборде и диаграмме Ганта, а также могут быть выгружены для анализа.
\item Обучающий модуль генерирует синтетические данные (правило LWKR) или использует исторические MILP-решения, тренирует модели с заданными гиперпараметрами.
\item Все изменения состояния ресурсов (готовность, ремонт) и заказов синхронизируются с базой данных в реальном времени.
\end{enumerate}

\paragraph{Функциональные требования к программной системе}

Функциональные требования к интеллектуальной системе производственного планирования определяют набор основных возможностей, необходимых для автоматизации диспетчеризации и повышения эффективности использования оборудования. Система должна обеспечивать ведение портфеля заказов, построение выполнимых расписаний, управление производственными ресурсами, обучение и применение нейросетевых моделей \cite{cifradvoynik}.

Функциональные требования:
\begin{enumerate}
\item Система должна предоставлять возможность создания, редактирования и удаления заказов с указанием номенклатуры, срока сдачи, приоритета и технологических операций.
\item Должна быть реализована функция построения производственного расписания (эвристикой или нейросетевым диспетчером) с назначением операций на ресурсы с учётом загрузки, готовности и специализации.
\item Необходимо обеспечить ручное управление очередями на дашборде: перенос задач между нераспределёнными и очередями ресурсов, запуск в производство, завершение операции, сброс всех задач в нераспределённые.
\item Система должна отображать актуальное состояние каждого ресурса: статус, готовность (надёжность), флаг ремонта, число задач в очереди, наличие активной задачи.
\item Пользователь должен иметь возможность изменять параметры ресурса (готовность, признак ремонта) с немедленным сохранением в базе данных.
\item Должна быть реализована возможность обучения нейросетевых моделей (PointerNet, SLIM, PPO) на сгенерированных датасетах с сохранением весов и метрик обучения в файловой системе.
\item Система должна позволять загружать и сохранять веса обученных моделей, а также автоматически подгружать последнюю использованную модель при старте.
\item Должна вестись история завершённых операций с возможностью просмотра и удаления записей.
\item Необходимо поддерживать управление технологическими типами операций и связями «ресурс–тип», ограничивающими назначение операций определённым набором станков.
\item Административные функции (добавление/удаление ресурсов и типов операций, настройка связей) доступны тому же пользователю без разделения ролей.
\end{enumerate}

В системе предусмотрено одно действующее лицо:
Пользователь – управляет заказами, ресурсами, очередями, обучением моделей и просмотром метрик, т.е. совмещает функции оператора и администратора.

На рисунке ~\ref{actorusing:image} представлена диаграмма прецедентов пользователя, содержащая функции, которые должны быть реализованы в разрабатываемой ИС.

\begin{figure}[H]
\centering
\includegraphics[height=0.8\textheight, width=\linewidth, keepaspectratio]{actorusing}
\caption{Диаграмма прецедентов пользователя}
\label{actorusing:image}
\end{figure}

\textbf{Прецедент <<Создать заказ>>}

Предусловия: Пользователь находится на вкладке «Заказы». В системе имеется хотя бы один тип операции.

Входные данные: название заказа, номер (опционально), дата и время сдачи, приоритет, тип операции, длительность, изделие, количество деталей.

Выходные данные: новый заказ с одной операцией, сохранённый в базе данных и отображённый на дашборде и в таблице заказов.

Постусловия: заказ доступен для планирования.

Сценарий:
\begin{enumerate}
\item Пользователь переходит на вкладку «Заказы».
\item Заполняет поля формы, выбирает тип операции из выпадающего списка, вводит длительность, изделие и количество.
\item Нажимает «Сохранить заказ».
\item Система проверяет корректность данных, создаёт заказ и операцию, присваивает уникальный идентификатор.
\item Заказ записывается в базу данных, появляется в списке заказов и в нераспределённых задачах на дашборде.
\end{enumerate}

\textbf{Прецедент <<Построить план>>}

Предусловия: в системе есть неразмещённые заказы.

Входные данные: нажатие кнопки «Построить план (RL)» или «Точный MILP-план».

Выходные данные: расписание, в котором каждой операции назначены ресурс, время начала и окончания; обновлённая диаграмма Ганта.

Постусловия: задачи помещены в очереди ресурсов, дашборд обновлён.

Сценарий:

\begin{enumerate}
\item Пользователь нажимает кнопку построения плана на дашборде или вкладке «Планирование».
\item Если загружена нейросеть, система формирует векторы признаков для каждой операции, пропускает через модель, корректирует логиты штрафом за низкую готовность и выбирает ресурс. Иначе применяется эвристика EDD и минимальная эффективная загрузка.
\item Полученное расписание записывается в current\_schedule и синхронизируется с таблицей task\_queue.
\item Задачи автоматически запускаются на исправных ресурсах (функция «В работу»).
\item Обновляются дашборд и диаграмма Ганта.
\end{enumerate}

\textbf{Прецедент <<Завершить операцию>>}

Предусловия: на ресурсе есть активная задача (статус active).

Входные данные: нажатие кнопки «Завершить» на карточке задачи.

Выходные данные: операция переведена в статус completed, соответствующий заказ удалён из активных, следующая ожидающая задача на ресурсе запущена.

Постусловия: очередь ресурса продвинута, запись о завершении добавлена в логи и вкладку «Завершённые».

Сценарий:

\begin{enumerate}
\item Пользователь на дашборде нажимает «Завершить» на активной задаче.
\item Система определяет ресурс, обновляет статус задачи, логирует событие с указанием ресурса.
\item Заказ, к которому относилась операция, помечается завершённым в БД и убирается из активного списка.
\item Система проверяет очередь освободившегося ресурса и, если он исправен, запускает первую ожидающую задачу.
\item Дашборд и вкладка «Завершённые» обновляются.
\end{enumerate}

\textbf{Прецедент <<Изменить готовность / ремонт ресурса>>}

Предусловия: ресурс существует и отображается на дашборде.

Входные данные: двойной клик по карточке ресурса, ввод нового значения готовности (0.0–1.0) и/или установка флажка «Ресурс в ремонте».

Выходные данные: обновлённые атрибуты ресурса, сохранённые в БД; перерисовка дашборда.

Постусловия: при следующем планировании изменённые параметры влияют на назначение операций.

Сценарий:

\begin{enumerate}
\item Пользователь дважды кликает по карточке ресурса.
\item Открывается диалог с полями «Готовность» и «Ремонт», а также информацией о текущей очереди.
\item Пользователь изменяет значение и/или флажок, нажимает «Сохранить».
\item Система записывает новые значения в объект ресурса и в базу данных, немедленно обновляя индикаторы на дашборде.
\end{enumerate}

\textbf{Прецедент <<Обучить нейросетевую модель>>}

Предусловия: имеется файл обучающей выборки gpss\_dataset.npz, пользователь находится на вкладке «Обучение».

Входные данные: количество эпох (или эпизодов), скорость обучения, выбор типа модели (PointerNet/SLIM/PPO).

Выходные данные: обученная модель (файл .pth), набор метрик в папке metrics/.

Постусловия: модель может быть загружена и использована для планирования.

Сценарий:

\begin{enumerate}
\item Пользователь переходит на вкладку «Обучение», при необходимости загружает датасет.
\item Указывает гиперпараметры и нажимает соответствующую кнопку (например, «Обучить PointerNet»).
\item Система загружает данные, создаёт модель SimpleDispatcher, запускает обучение, отображая прогресс в реальном времени.
\item По окончании сохраняет веса модели (pointer\_net.pth) и все метрики (кривые потерь, точности, матрицу ошибок) в папку metrics/<Модель>/<дата\_время>/.
\item Пользователь получает уведомление о завершении.
\end{enumerate}

\textbf{Прецедент <<Добавить ресурс>>}

Предусловия: пользователь находится на вкладке «Оборудование» → «Ресурсы».

Входные данные: идентификатор и название ресурса.

Выходные данные: новый ресурс, добавленный в БД и отображаемый на дашборде.

Постусловия: ресурс готов к планированию (готовность 1.0, не в ремонте).

Сценарий:

\begin{enumerate}
\item Пользователь вводит ID и название ресурса в соответствующие поля.
\item Нажимает «Добавить».
\item Система проверяет уникальность ID, создаёт объект Resource, сохраняет его в БД.
\item Ресурс появляется в списке ресурсов и на дашборде.
\end{enumerate}

Остальные прецеденты (редактирование заказа, удаление заказа, настройка связей, просмотр метрик и т.д.) реализованы аналогично и доступны через соответствующие вкладки интерфейса.

\paragraph{Требования пользователя к интерфейсу программной системы}

Пользовательский интерфейс программной системы должен быть интуитивно понятным и обеспечивать доступ к основным функциям без необходимости в специальных технических знаниях. Интерфейс предназначен для взаимодействия как с операторами, управляющими производственными заказами и расписаниями, так и с технологами, настраивающими оборудование и нейросетевые модели \cite{interface}.

Для пользователя интерфейс должен предоставлять возможность:
\begin{itemize}
\item создания, редактирования и удаления заказов на производство с указанием типа операции, длительности, количества и срока сдачи;
\item управления производственными ресурсами: изменения готовности (надёжности) и признака ремонта, просмотра текущей загрузки;
\item запуска построения плана с выбором метода (эвристический или нейросетевой) и визуализации полученного расписания на диаграмме Ганта;
\item ручного управления очередями: перенос задач из нераспределённых в очереди конкретных ресурсов, запуск задачи в производство, завершение операции, сброс всех задач обратно в нераспределённые;
\item просмотра истории завершённых операций и их удаления;
\item обучения нейросетевых моделей (PointerNet, SLIM, PPO) с настройкой гиперпараметров (число эпох/эпизодов, скорость обучения), отображением прогресса обучения и сохранением метрик;
\item загрузки и сохранения весов обученных моделей, а также автоматической подгрузки последней использованной модели при старте;
\item управления типами операций и связями «ресурс–тип», ограничивающими назначение операций.
\end{itemize}

Все вышеперечисленные функции доступны в рамках единого графического интерфейса, разделённого на тематические вкладки: «Дашборд», «Заказы», «Планирование», «Оборудование», «Состояния», «Обучение», «Настройки», «Завершённые». Основное взаимодействие с очередями и ресурсами осуществляется через дашборд, поддерживающий механизм drag-and-drop для перемещения задач.

Макет пользовательского интерфейса предоставлен на рисунке ~\ref{maketinterface:image}.

\begin{figure}[H]
\includegraphics[width=1\linewidth]{maketinterface}
\caption{Макет пользовательского интерфейса}
\label{maketinterface:image}
\end{figure}
\vspace{-14pt}
\paragraph{Нефункциональные требования к программной системе}

Нефункциональные требования к интеллектуальной системе производственного планирования определяют характеристики, обеспечивающие её надёжность, производительность и удобство эксплуатации.

Система должна обеспечивать высокую производительность при построении расписаний. Время формирования плана для типового объёма заказов (до 100 операций на 5–6 ресурсах) не должно превышать нескольких секунд как для эвристического алгоритма, так и для нейросетевого диспетчера.

Система должна быть отказоустойчивой и стабильной. Все изменения состояния ресурсов и заказов должны немедленно фиксироваться в базе данных, исключая потерю данных при нештатном завершении приложения. При возникновении ошибок во время обучения модели должно быть предусмотрено корректное завершение процесса с сохранением лога ошибок.

Обеспечение целостности данных является критически важным. База данных SQLite должна поддерживать ссылочную целостность между заказами, операциями, ресурсами и очередями. Все операции с данными выполняются в рамках одного приложения, поэтому дополнительная аутентификация не требуется, однако должна быть исключена возможность одновременной порчи данных из-за параллельных процессов.

Система должна быть масштабируемой в части увеличения числа заказов, ресурсов и типов операций без модификации исходного кода. Добавление новых ресурсов или типов операций должно выполняться через штатный интерфейс. Модели машинного обучения должны допускать повторное обучение на обновлённых выборках.

\paragraph{Требования к аппаратным средствам}
Для корректного функционирования интеллектуальной системы производственного планирования достаточно одного локального компьютера, выполняющего одновременно функции сервера (хранение данных, обучение моделей) и клиентского рабочего места (пользовательский интерфейс).

Компьютер должен быть оснащён многоядерным процессором с тактовой частотой не менее 1.5 ГГц, иметь объём оперативной памяти не менее 4 ГБ и экран с разрешением не менее 1280×800 пикселей. Для хранения базы данных, исполняемого кода, обученных моделей и метрик требуется не менее 500 МБ свободного дискового пространства.

Графический процессор (GPU) не является обязательным, однако его наличие может существенно ускорить обучение нейросетевых моделей при большом объёме данных.

\paragraph{Требования к программным средствам}

Программное обеспечение должно функционировать под управлением операционных систем семейства Windows (версия 10 или новее), Linux (основные дистрибутивы).

На компьютере должен быть установлен интерпретатор Python версии 3.9 или выше со следующими библиотеками:

\begin{itemize}
\item PySide6 (графический интерфейс);
\item PyTorch (обучение и применение нейросетевых моделей);
\item NumPy, Matplotlib, scikit-learn (обработка данных, визуализация метрик, расчёт метрик классификации);
\item SQLite3 (встроен в стандартную библиотеку Python, используется для хранения данных).
\end{itemize}

Дополнительное серверное программное обеспечение не требуется: система использует локальную базу данных SQLite, доступ к которой осуществляется непосредственно из приложения.

\paragraph{Требования к архитектуре нейронной сети}

Для решения задачи выбора оптимального ресурса при планировании используется искусственная нейронная сеть прямого распространения (многослойный перцептрон). Модель принимает на вход вектор признаков фиксированной размерности, описывающий одну технологическую операцию и текущее состояние всех доступных ресурсов \cite{keras}.

Входной вектор состоит из 36 признаков: 6 признаков операции (one-hot кодировка типа операции, нормализованная длительность, приоритет заказа, slack до срока сдачи) и для каждого из 5 ресурсов по 6 признаков (признак готовности, текущая загрузка, совместимость с типом операции, средняя историческая загрузка, флаг ремонта, коэффициент надёжности).

Архитектура модели включает три полносвязных слоя:
\begin{enumerate}
\item входной слой размерностью 36 нейронов;
\item скрытый слой из 64 нейронов с функцией активации ReLU и пакетной нормализацией (BatchNorm);
\item скрытый слой из 32 нейронов с функцией активации ReLU и пакетной нормализацией;
\item выходной слой из N нейронов (по числу ресурсов), выдающий ненормированные логиты вероятностей выбора каждого ресурса.
\end{enumerate}

Для принятия решения к выходным логитам применяется поправка, пропорциональная логарифму эффективной готовности ресурса, что позволяет явно учесть надёжность и ремонт при выборе.

Обучение модели выполняется по следующим парадигмам:

\begin{enumerate}
\item Поведенческое клонирование (BC, PointerNet). Минимизация кросс-энтропии между предсказанным и эталонным распределением, полученным эвристическим правилом минимальной эффективной загрузки.
\item Самоконтролируемое обучение с саморазметкой (SLIM). Аналогичная модель с периодическим добавлением уверенных предсказаний на неразмеченных данных.
\item Обучение с подкреплением (PPO). Модель обучается максимизировать отрицательное опоздание заказов, взаимодействуя с имитационной средой производственного планирования. 
\end{enumerate}

\subsubsection{Требования к оформлению документации}

Оформление документации должно осуществляться с учётом стандартов СТУ 02.030-2023, единой системы программной документации, единой системы конструкторской документации.

\subsection{Технический проект}                      %              Тех проект
\subsubsection{Общие сведения об интеллектуальной системе оптимизации}

Интеллектуальная система производственного планирования (APS‑система) предназначена для автоматизации процессов составления выполнимых расписаний и оперативной диспетчеризации на предприятиях с дискретным характером производства. Система предоставляет оператору возможность гибкого управления заказами, ресурсами и очередями, а также позволяет обучать и применять нейросетевые диспетчеры для повышения эффективности использования оборудования.

Программный продукт реализован в виде настольного приложения на языке Python с графическим интерфейсом PySide6 и локальной реляционной базой данных SQLite. Такая архитектура обеспечивает простоту развёртывания, автономность работы и низкие требования к инфраструктуре. Вся информация о заказах, операциях, ресурсах, технологических связях и истории исполнения хранится централизованно в единой базе данных, что гарантирует целостность и непротиворечивость данных \cite{isakova}.

Система спроектирована с учётом возможности расширения числа производственных ресурсов и типов операций без модификации исходного кода. Добавление нового станка или технологической операции выполняется через штатный пользовательский интерфейс, после чего они сразу становятся доступными для планирования. Модульная архитектура позволяет также подключать новые нейросетевые модели и алгоритмы планирования, что делает систему гибкой платформой для исследований и практического применения.

\subsubsection{Проект данных интеллектуальной системы}

При проектировании системы производственного планирования была использована локальная реляционная база данных SQLite, обеспечивающая хранение всех сущностей предметной области и их взаимосвязей. Структура базы данных разработана с учётом требований нормализации и поддержки ссылочной целостности.

Таблица «Заказы» содержит информацию о производственных заказах. Для каждого заказа хранится уникальный идентификатор, срок сдачи, приоритет, наименование и статус (активен или завершён). В таблице \ref{atr_zakaz:table} представлены атрибуты таблицы «Заказы».

\renewcommand{\arraystretch}{0.8} % уменьшение расстояний до сетки таблицы
\begin{xltabular}{\textwidth}{|p{3.5cm}|>{\setlength{\baselineskip}{0.7\baselineskip}}p{2.85cm}|X|}
\caption{Атрибуты таблицы <<Заказы>> \label{atr_zakaz:table}}\\
\hline \centrow \setlength{\baselineskip}{0.7\baselineskip} Атрибут & \centrow \setlength{\baselineskip}{0.7\baselineskip} Тип & \centrow Описание \\
\hline \centrow 1 & \centrow 2 & \centrow 3\\ 
\hline
\endfirsthead
\continuecaption{Продолжение таблицы \ref{atr_zakaz:table}}\\
\hline \centrow 1 & \centrow 2 & \centrow 3\\ 
\hline
\finishhead
id & string &  Уникальный идентификатор заказа (может быть с символами)\\
\hline due\_date & datetime & Плановая дата и время сдачи заказа\\
\hline priority\_weight & real & Весовой коэффициент приоритета (от 0.5 до 2.0)\\
\hline name & string & Краткое наименование заказа\\
\hline status & string & Текущее состояние заказа
\end{xltabular}
\renewcommand{\arraystretch}{1.0} % восстановление сетки

Таблица «Операции»  содержит детализацию технологических операций, входящих в заказы. Каждая операция связана с родительским заказом и характеризуется типом, нормативной длительностью, количеством изготавливаемых деталей и другими параметрами. В таблице \ref{atr_oper:table} приведены атрибуты таблицы «Операции».

\renewcommand{\arraystretch}{0.8} % уменьшение расстояний до сетки таблицы
\begin{xltabular}{\textwidth}{|p{3.5cm}|>{\setlength{\baselineskip}{0.7\baselineskip}}p{2.85cm}|X|}
\caption{Атрибуты таблицы <<Операции>> \label{atr_oper:table}}\\
\hline \centrow \setlength{\baselineskip}{0.7\baselineskip} Атрибут & \centrow \setlength{\baselineskip}{0.7\baselineskip} Тип & \centrow Описание \\
\hline \centrow 1 & \centrow 2 & \centrow 3\\ 
\hline
\endfirsthead
\continuecaption{Продолжение таблицы \ref{atr_oper:table}}\\
\hline \centrow 1 & \centrow 2 & \centrow 3\\ 
\hline
\finishhead
id & string & Уникальный идентификатор операции\\
\hline order\_id & string & Идентификатор заказа, к которому относится операция (внешний ключ)\\
\hline item & string & Наименование изготавливаемого изделия\\
\hline op\_number & integer & Порядковый номер операции в заказе\\
\hline norm\_duration & real & Нормативная длительность операции в минутах\\
\hline type\_id & string & Идентификатор типа операции \\
\hline predecessors & string & Список предшествующих операций (для реализации технологической последовательности)\\
\hline quantity & integer & Количество деталей\\
\hline urgency & real & Коэффициент срочности\\
\hline slack\_hours & real & Запас времени до дедлайна
\end{xltabular}
\renewcommand{\arraystretch}{1.0} % восстановление сетки

Таблица «Ресурсы»  описывает производственное оборудование. Для каждого ресурса задаются его наименование, текущий статус, нагрузочные характеристики и параметры готовности. В таблице \ref{atr_res:table} представлены атрибуты таблицы «Ресурсы».

\renewcommand{\arraystretch}{0.8} % уменьшение расстояний до сетки таблицы
\begin{xltabular}{\textwidth}{|p{3.9cm}|>{\setlength{\baselineskip}{0.7\baselineskip}}p{2.85cm}|X|}
\caption{Атрибуты таблицы <<Ресурсы>> \label{atr_res:table}}\\
\hline \centrow \setlength{\baselineskip}{0.7\baselineskip} Атрибут & \centrow \setlength{\baselineskip}{0.7\baselineskip} Тип & \centrow Описание \\
\hline \centrow 1 & \centrow 2 & \centrow 3\\ 
\hline
\endfirsthead
\continuecaption{Продолжение таблицы \ref{atr_res:table}}\\
\hline \centrow 1 & \centrow 2 & \centrow 3\\ 
\hline
\finishhead
id & string & Уникальный идентификатор ресурса\\
\hline name & string & Наименование ресурса \\
\hline section & string & Производственный участок \\
\hline status & string & Текущий статус \\
\hline operator\_name & string & ФИО оператора, закреплённого за ресурсом \\
\hline load\_minutes & real & Накопленная загрузка в минутах \\
\hline downtime\_minutes & real & Время простоя в минутах \\
\hline work\_hours & integer & Доступный фонд рабочего времени в часах \\
\hline repair & integer & Признак нахождения в ремонте\\
\hline reliability & real & Коэффициент готовности\\
\end{xltabular}
\renewcommand{\arraystretch}{1.0} % восстановление сетки

Таблица «Типы операций» содержит справочник технологических типов (фрезеровка, сборка, сварка и др.). В таблице ~\ref{atr_op_types:table} приведены её атрибуты.

\renewcommand{\arraystretch}{0.8} % уменьшение расстояний до сетки таблицы
\begin{xltabular}{\textwidth}{|p{4.5cm}|>{\setlength{\baselineskip}{0.7\baselineskip}}p{2.85cm}|X|}
\caption{Атрибуты таблицы <<Типы операций>> \label{atr_op_types:table}}\\
\hline \centrow \setlength{\baselineskip}{0.7\baselineskip} Атрибут & \centrow \setlength{\baselineskip}{0.7\baselineskip} Тип & \centrow Описание \\
\hline \centrow 1 & \centrow 2 & \centrow 3\\ 
\hline
\endfirsthead
\continuecaption{Продолжение таблицы \ref{atr_op_types:table}}\\
\hline \centrow 1 & \centrow 2 & \centrow 3\\ 
\hline
\finishhead
id & string & Уникальный идентификатор типа\\
\hline name & string & Название типа операции\\
\end{xltabular}
\renewcommand{\arraystretch}{1.0} % восстановление сетки

Таблица «Связи ресурсов с типами»  определяет, какие типы операций может выполнять каждый ресурс. Используется для назначения операций на подходящее оборудование. В таблице ~\ref{atr_res_op_types:table} перечислены атрибуты.

\renewcommand{\arraystretch}{0.8} % уменьшение расстояний до сетки таблицы
\begin{xltabular}{\textwidth}{|p{3.5cm}|>{\setlength{\baselineskip}{0.7\baselineskip}}p{2.85cm}|X|}
\caption{Атрибуты таблицы <<Связи ресурсов с типами>> \label{atr_res_op_types:table}}\\
\hline \centrow \setlength{\baselineskip}{0.7\baselineskip} Атрибут & \centrow \setlength{\baselineskip}{0.7\baselineskip} Тип & \centrow Описание \\
\hline \centrow 1 & \centrow 2 & \centrow 3\\ 
\hline
\endfirsthead
\continuecaption{Продолжение таблицы \ref{atr_res_op_types:table}}\\
\hline \centrow 1 & \centrow 2 & \centrow 3\\ 
\hline
\finishhead
resource\_id & string & Идентификатор ресурса \\
\hline operation\_type\_id & string & Идентификатор типа операции\\
\end{xltabular}
\renewcommand{\arraystretch}{1.0} % восстановление сетки

Таблица «Очередь задач» хранит текущее состояние очередей на ресурсах. Каждая запись связывает операцию с конкретным ресурсом, определяет её позицию в очереди и статус выполнения. В таблице ~\ref{atr_task_queue:table} даны атрибуты.

\renewcommand{\arraystretch}{0.8} % уменьшение расстояний до сетки таблицы
\begin{xltabular}{\textwidth}{|p{3.5cm}|>{\setlength{\baselineskip}{0.7\baselineskip}}p{2.85cm}|X|}
\caption{Атрибуты таблицы <<Очередь задач>> \label{atr_task_queue:table}}\\
\hline \centrow \setlength{\baselineskip}{0.7\baselineskip} Атрибут & \centrow \setlength{\baselineskip}{0.7\baselineskip} Тип & \centrow Описание \\
\hline \centrow 1 & \centrow 2 & \centrow 3\\ 
\hline
\endfirsthead
\continuecaption{Продолжение таблицы \ref{atr_task_queue:table}}\\
\hline \centrow 1 & \centrow 2 & \centrow 3\\ 
\hline
\finishhead
resource\_id & string & Идентификатор ресурса, в очередь которого помещена операция\\
\hline operation\_id & string & Идентификатор операции\\
\hline position & integer & Позиция в очереди\\
\hline status & string & Статус выполнения\\
\end{xltabular}
\renewcommand{\arraystretch}{1.0} % восстановление сетки

Таблица «Журнал состояний операций» фиксирует историю изменения статусов операций. Это необходимо для отслеживания хода производства и формирования отчётов. В таблице ~\ref{atr_op_state_log:table} представлены её атрибуты.

\renewcommand{\arraystretch}{0.8} % уменьшение расстояний до сетки таблицы
\begin{xltabular}{\textwidth}{|p{3.5cm}|>{\setlength{\baselineskip}{0.7\baselineskip}}p{2.85cm}|X|}
\caption{Атрибуты таблицы <<Журнал состояний операций>> \label{atr_op_state_log:table}}\\
\hline \centrow \setlength{\baselineskip}{0.7\baselineskip} Атрибут & \centrow \setlength{\baselineskip}{0.7\baselineskip} Тип & \centrow Описание \\
\hline \centrow 1 & \centrow 2 & \centrow 3\\ 
\hline
\endfirsthead
\continuecaption{Продолжение таблицы \ref{atr_op_state_log:table}}\\
\hline \centrow 1 & \centrow 2 & \centrow 3\\ 
\hline
\finishhead
id & integer & Автоинкрементный первичный ключ\\
\hline operation\_id & string & Идентификатор операции\\
\hline status & string & Статус операции на момент записи\\
\hline resource\_id & string & Идентификатор ресурса, на котором выполнялась операция\\
\hline state\_id & string & Идентификатор типа состояния \\
\hline timestamp & datetime & Метка времени события \\
\hline source & string & Источник изменения\\
\end{xltabular}
\renewcommand{\arraystretch}{1.0} % восстановление сетки

Таблица «Настройки приложения» хранит конфигурационные параметры в виде пар «ключ‑значение». Используется для сохранения путей к файлам моделей, гиперпараметров обучения и других пользовательских предпочтений. В таблице ~\ref{atr_app_settings:table} указаны атрибуты.

\renewcommand{\arraystretch}{0.8} % уменьшение расстояний до сетки таблицы
\begin{xltabular}{\textwidth}{|p{3.5cm}|>{\setlength{\baselineskip}{0.7\baselineskip}}p{2.85cm}|X|}
\caption{Атрибуты таблицы <<Настройки приложения>> \label{atr_app_settings:table}}\\
\hline \centrow \setlength{\baselineskip}{0.7\baselineskip} Атрибут & \centrow \setlength{\baselineskip}{0.7\baselineskip} Тип & \centrow Описание \\
\hline \centrow 1 & \centrow 2 & \centrow 3\\ 
\hline
\endfirsthead
\continuecaption{Продолжение таблицы \ref{atr_app_settings:table}}\\
\hline \centrow 1 & \centrow 2 & \centrow 3\\ 
\hline
\finishhead
key & string & Уникальное имя параметра\\
\hline value & string & Значение параметра\\
\hline updated\_at & datetime & Дата и время последнего обновления\\
\end{xltabular}
\renewcommand{\arraystretch}{1.0} % восстановление сетки

На рисунке ~\ref{db:image} представлена схема базы данных, отражающая основные таблицы и связи между ними.

\begin{figure}[H]
\includegraphics[width=1\linewidth]{db}
\caption{Структурная схема локальной базы данных}
\label{db:image}
\end{figure}

\subsubsection{Проектирование архитектуры программы}
\paragraph{Компоненты программы}

Разрабатываемая интеллектуальная система производственного планирования состоит из трёх логических групп компонентов, функционирующих в рамках одного настольного приложения.

Хранение информации обеспечивается реляционной базой данных SQLite, в которой размещены все основные сущности: заказы, операции, ресурсы, очереди задач, журналы состояний и настройки. Дополнительно используется файловое хранилище для размещения генерируемых обучающих выборок, обученных весов нейросетевых моделей и результатов метрик.

Для реализации системы используется язык программирования Python версии 3.9 и выше. Графический интерфейс построен с применением библиотеки PySide6, архитектура интерфейса основана на императивном подходе с выделением классов для каждой вкладки и диалогового окна. Бизнес-логика сосредоточена в классе APSCore, который координирует работу всех подсистем.

В процессе разработки используются следующие библиотеки:

\begin{enumerate}
\item PySide6 – библиотека для создания кроссплатформенных графических интерфейсов на языке Python, предоставляющая полный набор виджетов, поддержку тем оформления, работу с диаграммами и графическими сценами.
\item PyTorch – библиотека машинного обучения, используемая для построения, обучения и применения нейросетевых моделей диспетчеризации. Обеспечивает автоматическое дифференцирование, оптимизаторы и функции потерь.
\item NumPy – библиотека для работы с многомерными массивами, применяется при формировании векторов признаков операций и ресурсов, а также при расчёте метрик.
\item Matplotlib – библиотека для визуализации данных, используется для построения графиков обучения (кривые потерь и точности) и результатов сравнения методов.
\item scikit-learn – библиотека машинного обучения, применяется для расчёта метрик классификации (матрица ошибок, precision, recall, f1-score) при оценке качества нейросетевых диспетчеров.
\item SQLite3 – встроенная в Python библиотека для работы с реляционной базой данных SQLite, используется для хранения всех производственных данных и настроек.
\end{enumerate}

Компоненты графического интерфейса включают:

\begin{enumerate}
\item Главное окно – компонент, реализующий основное окно приложения с тематическими вкладками: Дашборд, Заказы, Планирование, Оборудование, Состояния, Обучение, Настройки, Завершённые.
\item Дашборд – компонент, отображающий текущее состояние ресурсов, очередей и нераспределённых задач. Поддерживает механизм drag-and-drop для перемещения задач между очередями и ресурсами.
\item Формы заказов и ресурсов – компоненты, обеспечивающие создание, редактирование и удаление заказов и ресурсов с сохранением изменений в базу данных.
\item Диаграмма Ганта – компонент на основе QGraphicsScene, визуализирующий построенное расписание с цветовой кодировкой по типам операций.
\item Панель обучения – компонент, предоставляющий элементы управления для запуска обучения моделей, настройки гиперпараметров и отображения прогресса в реальном времени.
\item Панель метрик – компонент, отображающий логи обучения и результаты сравнения методов в текстовом и графическом виде.
\end{enumerate}

Компоненты ядра системы включают:

\begin{enumerate}
\item APSCore – центральный компонент, реализующий бизнес-логику приложения: управление заказами и ресурсами, вызов методов планирования, синхронизацию расписаний с очередями, обработку завершения операций.
\item Модуль планирования – компонент, содержащий эвристический алгоритм EDD+LWKR с учётом готовности ресурсов, а также интерфейсы для вызова нейросетевых диспетчеров PointerNet, SLIM и PPO.
\item Модуль обучения – компонент, отвечающий за генерацию обучающих выборок, запуск тренировки моделей и сохранение весов и метрик.
\item Модуль сравнения – компонент, выполняющий тестирование обученных моделей на едином сценарии и вычисляющий ключевые производственные показатели.
\item Модуль управления данными – компонент, обеспечивающий взаимодействие с базой данных SQLite и файловым хранилищем.
\end{enumerate}
%Компоненты хранения данных включают:
%\begin{enumerate}
%\item База данных SQLite – реляционная база, содержащая таблицы заказов, операций, ресурсов, типов операций, связей, очередей задач, журналов состояний и настроек.
%\item Файловое хранилище – папки с генерируемыми датасетами, файлами весов обученных моделей и результатами метрик.
%\end{enumerate}

На рисунке ~\ref{components:image} представлена структурная схема программной системы, выполненная в виде блок-схемы UML.

\begin{figure}[ht]
\includegraphics[width=1\linewidth]{components}
\caption{Структурная схема компонентов программы}
\label{components:image}
\end{figure}

\paragraph{Архитектура интеллектуальной системы}

Интеллектуальная система производственного планирования представляет собой монолитное настольное приложение, построенное по трёхуровневой архитектуре в рамках одного процесса.

Уровни включают: слой представления (пользовательский интерфейс), слой бизнес-логики (ядро системы) и слой данных (база данных и файловое хранилище). Все компоненты функционируют на одном рабочем месте, взаимодействуя напрямую без сетевых протоколов.

Слой представления реализован на базе библиотеки PySide6 и содержит главное окно с тематическими вкладками, формы ввода заказов и параметров ресурсов, диаграмму Ганта, а также диалоговые окна для редактирования сущностей. Пользовательский интерфейс взаимодействует исключительно с ядром системы, отправляя запросы на выполнение операций и получая обновлённые данные для отображения.

Слой бизнес-логики представлен классом APSCore, который инкапсулирует всю прикладную логику: управление заказами и ресурсами, эвристическое и нейросетевое планирование, обучение моделей, сравнение методов и управление очередями. Ядро системы координирует работу вложенных модулей (модуль планирования, модуль обучения, модуль сравнения), которые вызываются по мере необходимости. Доступ к данным осуществляется через модуль управления данными, абстрагирующий работу с базой данных SQLite и файловой системой.

Слой данных включает реляционную базу данных SQLite, хранящую все основные сущности (заказы, операции, ресурсы, очереди, журналы состояний, настройки), и файловое хранилище, в котором размещаются генерируемые датасеты, обученные веса нейросетевых моделей, метрики обучения и результаты сравнения методов. База данных находится в локальном файле aps.db, что исключает необходимость развёртывания серверного программного обеспечения и упрощает перенос системы.

Взаимодействие между уровнями организовано следующим образом: графический интерфейс формирует запрос, ядро системы обрабатывает его, при необходимости обращаясь к модулям планирования или обучения, и через модуль управления данными сохраняет или извлекает информацию. Результаты возвращаются в интерфейс для визуализации.

На рисунке ~\ref{sysarch:image} представлена многоуровневая диаграмма архитектуры системы, отображающая перечисленные слои и их основные компоненты.

\begin{figure}[H]
\includegraphics[width=1\linewidth]{sysarch}
\caption{Диаграмма архитектуры системы}
\label{sysarch:image}
\end{figure}
\vspace{-16pt}
\paragraph{Проектирование пользовательского интерфейса}

На основе требований к пользовательскому интерфейсу, представленных в пункте технического задания, с помощью библиотеки PySide6 был разработан пользовательский интерфейс. Интерфейс построен по принципу тематических вкладок, каждая из которых предоставляет доступ к определённой функциональной группе. Главное окно приложения содержит следующие вкладки: Дашборд, Заказы, Планирование, Оборудование, Состояния, Обучение, Настройки, Завершённые.

На рисунке ~\ref{maket1:image} представлен макет интерфейса вкладки «Дашборд».

\begin{figure}[H]
\includegraphics[width=1\linewidth]{maket1}
\caption{Макет интерфейса вкладки <<Дашборд>>}
\label{maket1:image}
\end{figure}

На макете вкладки «Дашборд» представлено основное рабочее пространство оператора. В верхней части расположены выпадающий список для выбора даты планирования и кнопки управления: «Утвердить текущий план», «RL-совет», «В работу» и «Сбросить в нераспределённые». Левая часть вкладки отведена под нераспределённые задачи, которые отображаются в виде цветных карточек с информацией о заказе, типе операции и сроках. Правая часть представляет сетку, в которой для каждого ресурса показаны три колонки: название и статус ресурса, текущая выполняемая операция и очередь ожидающих задач. Карточки задач можно перетаскивать между нераспределёнными, очередью ресурса и обратно. Двойной клик по карточке ресурса открывает окно редактирования параметров ресурса.

На рисунке ~\ref{maket2:image} представлен макет интерфейса вкладки «Заказы».

\begin{figure}[H]
\includegraphics[width=1\linewidth]{maket2}
\caption{Макет интерфейса вкладки <<Заказы>>}
\label{maket2:image}
\end{figure}

На макете вкладки «Заказы» представлена форма управления производственными заказами. В верхней части расположена форма создания нового заказа с полями: название заказа, номер, дата и время сдачи, приоритет, тип операции, длительность, изделие и количество деталей. Ниже размещена кнопка «Сохранить заказ», а также кнопки для генерации десяти случайных заказов и загрузки заказов из ERP-файла. Под формой отображается таблица существующих заказов с колонками: идентификатор заказа, название, срок сдачи, количество операций, приоритет и кнопка удаления. Двойной клик по карточке заказа в любой части интерфейса открывает окно редактирования параметров заказа.

На рисунке ~\ref{maket3:image} представлен макет интерфейса вкладки «Планирование».

\begin{figure}[H]
\includegraphics[width=1\linewidth]{maket3}
\caption{Макет интерфейса вкладки <<Планирование>>}
\label{maket3:image}
\end{figure}

На макете вкладки «Планирование» расположены кнопки управления построением расписания: «Построить план (RL)», «Точный MILP-план (медленно)» и «Применить фиксированные и перестроить». Нижнюю часть вкладки занимает диаграмма Ганта, визуализирующая построенное расписание. На диаграмме операции отображаются в виде цветных блоков, привязанных к временной шкале и дорожкам ресурсов. Цвет блока соответствует типу операции.

На рисунке ~\ref{maket4:image} представлен макет интерфейса вкладки «Оборудование».

\begin{figure}[H]
\includegraphics[width=1\linewidth]{maket4}
\caption{Макет интерфейса вкладки <<Оборудование>>}
\label{maket4:image}
\end{figure}

На макете вкладки «Оборудование» представлены три подвкладки: Ресурсы, Типы и Связи. На подвкладке «Ресурсы» отображается список текущих ресурсов и форма для добавления нового ресурса с полями идентификатора и названия. На подвкладке «Типы» представлен список типов операций и форма для добавления нового типа. На подвкладке «Связи» расположены выпадающие списки для выбора ресурса и типа операции, кнопка «Добавить связь» и таблица существующих связей с возможностью удаления.

На рисунке ~\ref{maket5:image} представлен макет интерфейса вкладки «Обучение».

\begin{figure}[H]
\includegraphics[width=1\linewidth]{maket5}
\caption{Макет интерфейса вкладки <<Обучение>>}
\label{maket5:image}
\end{figure}

На макете вкладки «Обучение» представлены две подвкладки: Текущая НС и Обучение НС. На подвкладке «Текущая НС» отображается название загруженной модели и кнопка «Загрузить модель» для выбора файла весов. На подвкладке «Обучение НС» размещены поля для настройки гиперпараметров (количество эпох, скорость обучения), кнопки запуска обучения для каждого из трёх методов и кнопка остановки обучения. В левой части отображается график изменения функции потерь в реальном времени, в правой – текстовый лог обучения.

На рисунке ~\ref{maket6:image} представлен макет интерфейса вкладки «Завершённые заказы».

\begin{figure}[H]
\includegraphics[width=1\linewidth]{maket6}
\caption{Макет интерфейса вкладки <<Завершённые заказы>>}
\label{maket6:image}
\end{figure}

На макете вкладки «Завершённые заказы» представлена таблица с историей завершённых операций. Таблица содержит колонки: идентификатор операции, тип операции, заказ, ресурс и кнопка удаления записи.

\subsection{Рабочий проект}                   %           3. раб проект

\subsubsection{Спецификация компонентов и классов интеллектуальной системы}

В приложении~А представлены фрагменты исходного кода разработанной интеллектуальной системы производственного планирования.

Класс APSCore реализует центральную бизнес-логику приложения: управление заказами и ресурсами, построение расписаний, запуск обучения и обработку очередей. В таблицах ~\ref{apscore_fields:table} и ~\ref{apscore_methods:table} приведены спецификации его основных полей и методов.

\begin{xltabular}{\textwidth}{|p{4.5cm}|>{\setlength{\baselineskip}{0.7\baselineskip}}p{2.85cm}|X|}
\caption{Поля класса APSCore \label{apscore_fields:table}}\\
\hline \centrow \setlength{\baselineskip}{0.7\baselineskip} Имя & \centrow Тип & \centrow Описание \\
\hline \centrow 1 & \centrow 2 & \centrow 3\\ 
\hline
\endfirsthead
\continuecaption{Продолжение таблицы \ref{apscore_fields:table}}\\
\hline \centrow Имя & \centrow Тип & \centrow Описание \\
\hline
\finishhead
resources & list[Resource] & Список производственных ресурсов\\
\hline resource\_map & dict[str, Resource] & Словарь ресурсов по идентификатору\\
\hline orders & list[Order] & Список активных заказов\\
\hline current\_schedule & dict[str, dict] & Текущее оперативное расписание\\
\hline approved\_schedule & dict[str, dict] & Утверждённое (фиксированное) расписание\\
\hline rl\_agent & SimpleAgent или None & Загруженный нейросетевой диспетчер\\
\hline settings & dict[str, str] & Словарь настроек приложения, загружаемый из БД\\
\hline cached\_res\_op\_types & dict[str, list] & Кэш связей ресурс-тип для ускорения планирования\\
\hline
\end{xltabular}
\vspace{-16pt}
\begin{xltabular}{\textwidth}{|p{7cm}|X|}
\caption{Основные методы класса APSCore \label{apscore_methods:table}}\\
\hline \centrow Имя & \centrow Описание \\
\hline \centrow 1 & \centrow 2\\ 
\hline
\endfirsthead
\continuecaption{Продолжение таблицы \ref{apscore_methods:table}}\\
\hline \centrow Имя & \centrow Описание \\
\hline
\finishhead
create\_order & Создаёт новый заказ с одной операцией и сохраняет в БД\\
\hline delete\_order & Удаляет заказ и все связанные с ним операции из БД и очередей\\
\hline generate\_random\_orders & Генерирует заданное количество случайных заказов\\
\hline plan\_with\_reliability & Строит расписание эвристикой EDD + эффективная загрузка\\
\hline run\_milp & Вызывает эвристическое планирование (заменяет точный MILP)\\
\hline \_sync\_schedule\_to\_task\_queue & Переносит назначения из расписания в очередь задач\\
\hline start\_all\_queues & Запускает первые ожидающие задачи на всех свободных ресурсах\\
\hline complete\_task & Завершает активную операцию, удаляет заказ и продвигает очередь\\
\hline load\_rl\_model\_if\_available & Загружает нейросетевую модель из файла, указанного в настройках\\
\hline plan\_with\_rl & Строит план с использованием загруженной нейросети\\
\hline
\end{xltabular}

Класс SimpleDispatcher реализует нейросетевую модель прямого распространения, предсказывающую наиболее подходящий ресурс для операции. Его спецификация приведена в таблицах ~\ref{simpledisp_fields:table} и ~\ref{simpledisp_methods:table}.

\begin{xltabular}{\textwidth}{|p{3.5cm}|>{\setlength{\baselineskip}{0.7\baselineskip}}p{2.85cm}|X|}
\caption{Поля класса SimpleDispatcher \label{simpledisp_fields:table}}\\
\hline \centrow Имя & \centrow Тип & \centrow Описание \\
\hline \centrow 1 & \centrow 2 & \centrow 3\\ 
\hline
\endfirsthead
\continuecaption{Продолжение таблицы \ref{simpledisp_fields:table}}\\
\hline \centrow Имя & \centrow Тип & \centrow Описание \\
\hline
\finishhead
net & nn.Sequential & Последовательная модель\\
\hline
\end{xltabular}

Последовательная модель нейросети в классе SimpleDispatcher содержит следующие слои: Linear(\(\rightarrow\)64)\(\rightarrow\)ReLU\(\rightarrow\)BatchNorm\(\rightarrow\)
Linear(\(\rightarrow\)32)\(\rightarrow\)ReLU\(\rightarrow\)BatchNorm\(\rightarrow\)Linear(\(\rightarrow\)N).

\begin{xltabular}{\textwidth}{|p{3cm}|X|}
\caption{Методы класса SimpleDispatcher \label{simpledisp_methods:table}}\\
\hline \centrow Имя & \centrow Описание \\
\hline \centrow 1 & \centrow 2\\ 
\hline
\endfirsthead
\continuecaption{Продолжение таблицы \ref{simpledisp_methods:table}}\\
\hline \centrow Имя & \centrow Описание \\
\hline
\finishhead
forward & Прямой проход: принимает тензор признаков (batch \(\times\) input\_dim) и возвращает логиты для каждого ресурса\\
\hline
\end{xltabular}

Класс Trainer обеспечивает обучение модели методом поведенческого клонирования (PointerNet). Спецификация приведена в таблицах ~\ref{trainer_fields:table} и ~\ref{trainer_methods:table}.

\begin{xltabular}{\textwidth}{|p{2.5cm}|>{\setlength{\baselineskip}{0.7\baselineskip}}p{3.8cm}|X|}
\caption{Поля класса Trainer \label{trainer_fields:table}}\\
\hline \centrow Имя & \centrow Тип & \centrow Описание \\
\hline \centrow 1 & \centrow 2 & \centrow 3\\ 
\hline
\endfirsthead
\continuecaption{Продолжение таблицы \ref{trainer_fields:table}}\\
\hline \centrow Имя & \centrow Тип & \centrow Описание \\
\hline
\finishhead
model & SimpleDispatcher & Обучаемая модель\\
\hline history & dict & Словарь с историей обучения (потери, точность)\\
\hline signals & TrainingSignals & Сигналы для отображения прогресса в GUI\\
\hline
\end{xltabular}
\vspace{-16pt}
\begin{xltabular}{\textwidth}{|p{7cm}|X|}
\caption{Методы класса Trainer \label{trainer_methods:table}}\\
\hline \centrow Имя & \centrow Описание \\
\hline \centrow 1 & \centrow 2\\ 
\hline
\endfirsthead
\continuecaption{Продолжение таблицы \ref{trainer_methods:table}}\\
\hline \centrow Имя & \centrow Описание \\
\hline
\finishhead
train\_bc & Загружает датасет, создаёт модель, запускает обучение и сохраняет метрики\\
\hline save & Сохраняет веса модели в файл\\
\hline load & Загружает веса модели из файла\\
\hline save\_metrics & Сохраняет графики обучения и отчёт о классификации в папку metrics\\
\hline
\end{xltabular}

Класс SelfLabelingTrainer (SLIM) расширяет обучение механизмом саморазметки. Его спецификация приведена в таблицах ~\ref{slimtrainer_fields:table} и ~\ref{slimtrainer_methods:table}.

\begin{xltabular}{\textwidth}{|p{3.5cm}|>{\setlength{\baselineskip}{0.7\baselineskip}}p{3.85cm}|X|}
\caption{Поля класса SelfLabelingTrainer \label{slimtrainer_fields:table}}\\
\hline \centrow Имя & \centrow Тип & \centrow Описание \\
\hline \centrow 1 & \centrow 2 & \centrow 3\\ 
\hline
\endfirsthead
\continuecaption{Продолжение таблицы \ref{slimtrainer_fields:table}}\\
\hline \centrow Имя & \centrow Тип & \centrow Описание \\
\hline
\finishhead
model & SimpleDispatcher & Обучаемая модель\\
\hline history & dict & Словарь с историей обучения\\
\hline signals & TrainingSignals & Сигналы для связи с GUI\\
\hline
\end{xltabular}
\vspace{-16pt}
\begin{xltabular}{\textwidth}{|p{7cm}|X|}
\caption{Методы класса SelfLabelingTrainer \label{slimtrainer_methods:table}}\\
\hline \centrow Имя & \centrow Описание \\
\hline \centrow 1 & \centrow 2\\ 
\hline
\endfirsthead
\continuecaption{Продолжение таблицы \ref{slimtrainer_methods:table}}\\
\hline \centrow Имя & \centrow Описание \\
\hline
\finishhead
train\_slim & Выполняет обучение с периодической саморазметкой неразмеченных данных\\
\hline save\_metrics & Сохраняет графики и отчёт в папку metrics\\
\hline
\end{xltabular}

Класс SimpleAgent служит обёрткой для использования обученной модели в процессе планирования. Спецификация приведена в таблицах ~\ref{simpleagent_fields:table} и ~\ref{simpleagent_methods:table}.

\begin{xltabular}{\textwidth}{|p{3.5cm}|>{\setlength{\baselineskip}{0.7\baselineskip}}p{3.85cm}|X|}
\caption{Поля класса SimpleAgent \label{simpleagent_fields:table}}\\
\hline \centrow Имя & \centrow Тип & \centrow Описание \\
\hline \centrow 1 & \centrow 2 & \centrow 3\\ 
\hline
\endfirsthead
\continuecaption{Продолжение таблицы \ref{simpleagent_fields:table}}\\
\hline \centrow Имя & \centrow Тип & \centrow Описание \\
\hline
\finishhead
model & SimpleDispatcher & Загруженная обученная модель\\
\hline device & torch.device & Устройство для вычислений (CPU)\\
\hline
\end{xltabular}

\begin{xltabular}{\textwidth}{|p{7cm}|X|}
\caption{Методы класса SimpleAgent \label{simpleagent_methods:table}}\\
\hline \centrow Имя & \centrow Описание \\
\hline \centrow 1 & \centrow 2\\ 
\hline
\endfirsthead
\continuecaption{Продолжение таблицы \ref{simpleagent_methods:table}}\\
\hline \centrow Имя & \centrow Описание \\
\hline
\finishhead
select\_action & Принимает вектор признаков (36) и возвращает индекс выбранного ресурса\\
\hline
\end{xltabular}

Также в разработанной программе имеются классы Resource, Order и Operation, которые являются датаклассами и не содержат специфических методов.

На рисунке ~\ref{classdiag:image} представлена диаграмма классов разработанной системы.

\begin{figure}[H]
\includegraphics[width=1\linewidth]{classdiag}
\caption{Диаграмма классов интеллектуальной системы }
\label{classdiag:image}
\end{figure}


\subsubsection{Системное тестирование разработанного приложения}

Для системного тестирования интеллектуальной системы производственного планирования были проведены проверки основных функциональных сценариев: управление заказами, построение расписаний, ручное управление очередями, завершение операций, изменение параметров ресурсов, обучение нейросетевых моделей и просмотр истории\cite[с.~109--125]{sidorkina}. Тестирование подтвердило, что программное обеспечение полностью соответствует заявленным функциональным требованиям.

На рисунке ~\ref{systest1:image} представлен пример создания нового заказа на вкладке «Заказы». Заполнены поля: название, срок сдачи, приоритет, тип операции, длительность, изделие и количество. После нажатия кнопки «Сохранить заказ» он появляется в таблице заказов и в нераспределённых задачах на дашборде.

\begin{figure}[H]
\includegraphics[width=1\linewidth]{systest1}
\caption{Пример создания нового заказа}
\label{systest1:image}
\end{figure}

На рисунке ~\ref{systest2:image} показан дашборд после добавления нескольких заказов. В левой части отображаются нераспределённые задачи (жёлтые карточки), в правой – ресурсы с индикаторами готовности, очереди задач и текущие активные операции.

\begin{figure}[H]
\includegraphics[width=1\linewidth]{systest2}
\caption{Дашборд с нераспределёнными задачами и состоянием ресурсов}
\label{systest2:image}
\end{figure}

На рисунке ~\ref{systest3:image} представлен результат построения плана (нажатием кнопки «Построить план (RL)»). Задачи распределены по ресурсам с учётом их готовности и загрузки. Активные операции отображаются зелёными карточками, ожидающие – синими.

\begin{figure}[H]
\includegraphics[width=1\linewidth]{systest3}
\caption{Дашборд после построения плана}
\label{systest3:image}
\end{figure}

На рисунке ~\ref{systest4:image} показана диаграмма Ганта на вкладке «Планирование». Цветные блоки соответствуют операциям разных типов, каждый блок привязан к своему ресурсу и временной шкале.

\begin{figure}[H]
\includegraphics[width=1\linewidth]{systest4}
\caption{Диаграмма Ганта построенного расписания}
\label{systest4:image}
\end{figure}

На рисунке ~\ref{systest5:image} продемонстрирована возможность ручного управления: задача перетащена из нераспределённых в очередь другого ресурса и сразу запущена в производство. Карточка задачи изменила цвет и положение.

\begin{figure}[H]
\includegraphics[width=1\linewidth]{systest5}
\caption{Ручное перемещение задачи и запуск в производство}
\label{systest5:image}
\end{figure}

На рисунке ~\ref{systest6:image} показан процесс завершения операции. После нажатия кнопки «Завершить» на активной задаче она исчезает из дашборда, а следующая ожидающая задача на этом ресурсе автоматически переходит в производство. Завершённый заказ появляется во вкладке «Завершённые заказы».

\begin{figure}[H]
\includegraphics[width=1\linewidth]{systest6}
\caption{Завершение операции и автоматическое продвижение очереди}
\label{systest6:image}
\end{figure}

На рисунке ~\ref{systest7:image} представлено окно редактирования параметров ресурса, вызываемое двойным кликом по карточке ресурса. Изменены значения готовности и ремонта. После сохранения индикаторы на дашборде обновились.

\begin{figure}[H]
\includegraphics[width=1\linewidth]{systest7}
\caption{Редактирование параметров ресурса}
\label{systest7:image}
\end{figure}

На рисунке ~\ref{systest8:image} показана вкладка «Обучение» во время тренировки нейросетевой модели. В левой части отображается график изменения функции потерь, в правой – текстовый лог с информацией об эпохах и точности.

\begin{figure}[H]
\includegraphics[width=1\linewidth]{systest8}
\caption{Процесс обучения нейросетевой модели}
\label{systest8:image}
\end{figure}

На рисунке ~\ref{systest9:image} представлена вкладка «Завершённые заказы» с таблицей выполненных операций. Для каждой записи указаны идентификатор операции, тип, заказ и ресурс, имеется кнопка удаления.

\begin{figure}[H]
\includegraphics[width=1\linewidth]{systest9}
\caption{Вкладка «Завершённые заказы» с историей операций}
\label{systest9:image}
\end{figure}

Проведённые тесты подтвердили корректную работу всех заявленных функций системы.

\subsubsection{Сборка компонентов интеллектуальной системы}

Сборка и запуск интеллектуальной системы производственного планирования выполняются непосредственно из исходного кода на языке Python, без предварительной компиляции. Для работы необходим интерпретатор Python версии 3.9 или выше и библиотеки, перечисленные в файле требований (PySide6, PyTorch, NumPy, Matplotlib, scikit-learn). Установка зависимостей производится однократно.

Графический интерфейс приложения реализован на библиотеке PySide6, бизнес-логика системы сосредоточена в ядре, построенном по модульному принципу. Модуль управления данными обеспечивает взаимодействие с реляционной базой данных SQLite. Модели предметной области описывают производственные ресурсы, заказы и технологические операции. Обучение нейросетевых моделей реализовано отдельными компонентами, поддерживающими поведенческое клонирование, самообучение с саморазметкой и обучение с подкреплением. Генерация обучающих выборок выполняется специализированным генератором, а сравнение методов планирования — аналитическим модулем, который вычисляет ключевые производственные показатели и строит сравнительные графики. Вспомогательные компоненты отвечают за управление путями к файлам весов, предсказание длительности операций и формирование тестовых производственных данных \cite{kachestvo2}.

Все исходные модули организованы в едином каталоге проекта, что обеспечивает простоту навигации и сопровождения. Дополнительные данные, такие как обученные веса, датасеты и метрики, размещаются в автоматически создаваемых подкаталогах. Запуск приложения выполняется одной командой, после чего открывается главное окно с тематическими вкладками, и система готова к работе.

